\documentclass[11pt]{article}

\usepackage[margin=1in]{geometry}
\usepackage{amsmath,amssymb}
\usepackage{microtype}
\usepackage[T1]{fontenc}
% Times-like text and math; falls back gracefully if newtx is unavailable.
\IfFileExists{newtxtext.sty}{\usepackage{newtxtext}\usepackage{newtxmath}}{\usepackage{lmodern}}
\usepackage[hidelinks]{hyperref}

\newcommand{\Shat}{\widehat{\Sigma}}
\newcommand{\tr}{\operatorname{tr}}
\newcommand{\logdet}{\log\det}

\title{\vspace{-2em}Schur Covariance Evaluation\\[0.3em]
       \large A Principled Pseudo-Likelihood in High Dimensions}
\author{Peter Cotton}
\date{Draft --- \texttt{microprediction/precise}}

\begin{document}
\maketitle

\begin{abstract}
\noindent
Ask a covariance matrix how surprised it is by fresh data and you get the Gaussian likelihood --- the
obvious, principled way to score an estimate. In high dimensions the obvious thing falls apart. The
likelihood lives or dies on the smallest eigenvalues of the estimate, which are exactly the ones you
cannot pin down when you have about as many variables as observations; so, as I will show, it ranks
estimators worse than a coin flip. This paper is about a one-knob repair. The Gaussian likelihood
factors, through Schur complements, into a sum of block-conditional pieces, and a single dial
$\gamma$ damps how much each block is allowed to trust the others. At $\gamma=1$ you have the full
likelihood; at $\gamma=0$ the block-diagonal (composite) likelihood that ignores the coupling
altogether; in between, something better conditioned than either. I call it the \emph{Schur
pseudo-likelihood}. It is the same Schur-complement dial that interpolates hierarchical risk parity
and the minimum-variance portfolio --- borrowed from allocation and pointed at evaluation instead. I
work out what $\gamma$ does exactly in a tractable case, where the best setting turns out to be the
\emph{reliability} of the coupling (a Wiener / James--Stein shrinkage); I show where it helps and
where it does not; and I try to be honest about which claims I can actually stand behind.
\end{abstract}

\section{The likelihood quietly fails}

If you want to know whether a covariance estimate is any good, the textbook answer is to score it by
the likelihood it assigns to held-out data. For a zero-mean Gaussian that is
\[
  \ell(\Shat;x) \;=\; -\tfrac12\bigl(\logdet\Shat + x^\top \Shat^{-1} x\bigr),
\]
or, averaging over a test sample and writing $S$ for the sample covariance,
$\ell(\Shat;S) = -\tfrac12\bigl(\logdet\Shat + \tr(\Shat^{-1}S)\bigr)$ --- the same thing the Wishart
likelihood and Stein's loss are measuring, dressed differently. It is a strictly proper scoring rule.
It is what the maximum-likelihood estimator is built to optimize. By every classical argument it is
the right answer, and in low dimensions it is.

So here is the awkward part. Both terms above are run by the \emph{smallest} eigenvalues of $\Shat$:
the log-determinant is $\sum_i \log\lambda_i$, which dives toward $-\infty$ as any $\lambda_i\to 0$,
and $\Shat^{-1}$ blows up in the same limit. When the dimension $p$ is comparable to the sample size
$n$ --- the regime everyone actually works in --- those small eigenvalues collapse toward zero and,
worse, become statistically unidentifiable: the Mar\v{c}enko--Pastur bulk is noise pretending to be
signal. The likelihood then stakes its entire verdict on the part of the spectrum it knows least
about. You can guess how that goes.

It is not subtle. In a controlled experiment where I know which of two estimates is genuinely better
(I can compute the true Kullback--Leibler divergence, because I built the truth), the held-out
likelihood picks the better one about $45\%$ of the time. A coin would do better. Inversion-free
judges --- Frobenius distance, the variogram score --- sail along at around $85\%$. The likelihood, the
one estimator we were told to trust, is the one that fails. This is the evaluation-side twin of
Markowitz's ``error maximization'': the minimum-variance portfolio $w \propto \Shat^{-1}\mathbf{1}$
falls apart for the very same reason, the very same $\Shat^{-1}$, the very same small eigenvalues. The
cure, in both places, is to stop trusting the raw inverse. The question is how to stop trusting it
\emph{by degrees}, with one interpretable knob, rather than throwing the coupling away entirely.

\section{What the Schur complement is doing here}

Nothing in this paper works without one exact fact, so let me state it carefully before I start
bending it. Split the variables into $K$ ordered blocks, $x=(x_1,\dots,x_K)$. The joint density
factors --- exactly, no approximation --- into a product of block conditionals,
$p(x)=\prod_k p(x_k\mid x_{<k})$, and for a Gaussian each conditional is itself Gaussian with
covariance the \emph{Schur complement}
\[
  S_k \;=\; \Sigma_{kk} - \Sigma_{k,<k}\,\Sigma_{<k,<k}^{-1}\,\Sigma_{<k,k}
\]
and mean $\mu_{k\mid<k} = \Sigma_{k,<k}\Sigma_{<k,<k}^{-1}x_{<k}$. So
$\logdet\Sigma = \sum_k \logdet S_k$, the quadratic form splits block by block, and the whole
likelihood is just a sum of block-conditional terms, each of which only ever inverts the conditioning
block $\Sigma_{<k,<k}$. This is the block Cholesky factorization, and it is the engine underneath the
Vecchia approximation, Gaussian--Markov-random-field likelihoods, and nested-dissection elimination.
I have made no approximation; I have only rewritten the likelihood in a form that lets me reach in and
turn down one thing.

\section{One dial}

Here is the dial. Damp the Schur complement, and the regression that comes with it, by a single number
$\gamma\in[0,1]$:
\[
  S_k(\gamma) \;=\; \Sigma_{kk} - \gamma\,\Sigma_{k,<k}\Sigma_{<k,<k}^{-1}\Sigma_{<k,k}
              \;=\; (1-\gamma)\,\Sigma_{kk} + \gamma\,S_k,
  \qquad
  \mu_{k\mid<k}(\gamma) = \gamma\,\Sigma_{k,<k}\Sigma_{<k,<k}^{-1}x_{<k},
\]
and score with the damped conditionals,
$\ell_\gamma(\Sigma;x) = \sum_k \log \mathcal{N}\!\bigl(x_k;\,\mu_{k\mid<k}(\gamma),\,S_k(\gamma)\bigr)$.
At $\gamma=1$ this is the exact joint likelihood again. At $\gamma=0$ each block is scored under its
own unconditional covariance $\Sigma_{kk}$, with no reference to the others --- a product of block
marginals, which is to say composite likelihood in the sense of Besag, Lindsay, and Varin--Reid--Firth.
Everywhere in between is the new object, and the only one that changes the block covariances rather
than the data: it inverts nothing larger than a conditioning block, and it is better conditioned than
the full likelihood precisely because damping lifts the small eigenvalues that sank it. I call it the
\emph{Schur pseudo-likelihood}.

Let me be honest about what it is, because the honesty is load-bearing later. At $\gamma=1$ it is a
normalized joint density. At $\gamma=0$ it is a composite likelihood: a sum of perfectly good
component log-densities that does not itself integrate to one. For $\gamma$ in between, each term is
still a genuine Gaussian log-density --- a strictly proper score on its own block --- but the damped
conditionals are not mutually consistent; they are not the conditionals of any single joint Gaussian.
So $\ell_\gamma$ is a \emph{pseudo}-likelihood, and that is the bargain: I trade the joint-density
property for components that stay numerically sane. As we will see in a moment, you must not try to
\emph{estimate} a covariance by maximizing it freely; it is a device for \emph{scoring} a covariance,
or for damping one you already have.

\section{What the dial actually does}

I can say exactly what $\gamma$ does, at least in the smallest case that still has anything to say.
Take two scalar blocks: $x_1\sim\mathcal{N}(0,a)$ and $x_2 = b\,x_1 + \varepsilon$ with
$\varepsilon\sim\mathcal{N}(0,s)$, so the true law of $x_2$ given $x_1$ has slope $b$ and conditional
variance $s$. Maximize the population $\ell_\gamma$ over candidate covariances and something
slightly unsettling happens. The maximizer recovers the true predictive law for \emph{every}
$\gamma$ --- the effective slope comes out to $b$ and the conditional variance to $s$ --- but the
covariance it implies is inflated, with off-diagonal $\widehat{\Sigma}_{12} = \Sigma_{12}^\star/\gamma$.
Push $\gamma$ down far enough and the implied matrix leaves the positive-definite cone altogether. The
damping, in other words, gets undone by rescaling: a free fit simply scales its coupling up by
$1/\gamma$ to compensate, and for strong enough coupling it scales right past a valid covariance. That
is the precise reason $\ell_\gamma$ is no good as a free estimation objective, and is fine as a score
or as a fixed transform --- both of which hold $\widehat{\Sigma}$ still instead of letting it absorb
the dial.

Now the part worth the trip. The thing $\gamma$ is really controlling is how much to trust the
estimated coupling, and \emph{that} has a clean optimum. Fit the slope from $n$ samples by least
squares; $\hat b$ is unbiased with variance $s/(a(n-2))$. Choosing $\gamma$ to minimize the expected
predictive residual gives
\[
  \boxed{\;\gamma^\star \;=\; \frac{b^2}{\,b^2 + \operatorname{Var}\hat b\,}
        \;=\; \frac{(n-2)\,\rho^2}{(n-2)\,\rho^2 + (1-\rho^2)}\;,\qquad
        \rho^2 = \frac{b^2 a}{\,b^2 a + s\,}.}
\]
Read it and the whole thesis is sitting there. $\gamma^\star$ is the \emph{reliability} of the
coupling --- a Wiener filter, a James--Stein shrinkage, whatever you like to call the signal-to-total
ratio. It climbs to $1$ as the sample grows or the coupling strengthens (trust it; use the full
likelihood) and falls to $0$ as the coupling dissolves into noise (ignore it; go block-diagonal), and
it sits in the interior exactly when the coupling is worth something but not everything. I checked it
against the simulated optimum and it lands; the vector version replaces $\rho^2$ with the largest
squared canonical correlation between the blocks, and across $K$ blocks the inflation, and the
positive-definite threshold, compound down the chain.

\section{Why an interior $\gamma$ wins, and where it doesn't}

It pays to be careful here, because the easy story is wrong. Write the score as a quadratic form,
$\ell_\gamma(\widehat{\Sigma};x) = -\tfrac12\bigl(\mathrm{Ld}_\gamma + x^\top W_\gamma(\widehat{\Sigma})\,x\bigr)$,
where $W_\gamma$ assembles the damped block conditionals and, at $\gamma=1$, is exactly
$\widehat{\Sigma}^{-1}$. For two fixed estimates and a Gaussian test sample the probability of ranking
the better one above the worse is governed by a per-sample signal-to-noise ratio, and that ratio --- I
worked it out and checked it against Monte Carlo --- is \emph{maximized at $\gamma=1$}. Under
test-sample noise alone, Neyman--Pearson wins and damping only adds bias. So the high-dimensional
advantage is not a test-noise effect, and any paper that claims otherwise is hand-waving.

Where it comes from is the instability of the \emph{estimate}. The score depends on $\widehat{\Sigma}$
through $W_\gamma$, and at $\gamma=1$, $W_1=\widehat{\Sigma}^{-1}$ has a sensitivity that scales like
$1/\lambda_{\min}^2$; the variance of the score gap across training draws then runs like
$1/\lambda_{\min}^4$, which I can watch explode --- from essentially nothing to dozens --- as the truth
gets more ill-conditioned, while a damped $\gamma$ keeps it bounded. When the discriminating signal
lives in the within-block structure and the cross-block coupling is shared noise, an interior $\gamma$
beats both ends outright (power around $0.84$ against $0.45$ at either extreme). The dial is a
bias--variance knob, but over the \emph{estimate}, not the test point.

One more honest line in the ledger. The estimating equation $\nabla_\Sigma \ell_\gamma$ is unbiased at
the truth only at the two ends --- $\gamma=1$ (the Fisher score) and $\gamma=0$ (correct block
marginals); in the interior it is biased, which is the inflation from two sections ago wearing a
different hat. So the classical Godambe efficiency story applies only at the endpoints, and the
interior is a \emph{tempering}, justified by the discrimination it buys rather than by being any kind
of consistent estimator. That is the rigorous form of ``score with it, don't fit with it.''

\section{It is a shrinkage, and a familiar one}

The well-behaved use --- damp the coupling of a covariance you already have --- is just shrinkage,
worth saying plainly. Each conditional covariance is pulled from its full Schur complement back toward
the unconditional block, $S_k(\gamma)=(1-\gamma)\Sigma_{kk}+\gamma S_k$, with intensity $1-\gamma$.
Since conditioning only ever shrinks variance, $S_k \preceq \Sigma_{kk}$, this raises the conditional
covariance and lifts the small eigenvalues of every block-sized inverse the score has to perform. It
is James--Stein applied blockwise to the conditioning. And it is orthogonal to the shrinkage everyone
already knows: Ledoit--Wolf and the rotation-invariant estimators clean the \emph{eigenvalues},
basis-free; this cleans the \emph{cross-block coupling}, basis-dependent. They compose. $\gamma$ is to
the block factorization what shrinkage intensity is to the spectrum.

\section{One knob, several familiar faces}

The reason I think this is worth a paper rather than a footnote is that the same dial keeps turning up,
and naming it organizes a surprising amount of scattered machinery onto two axes of one factorization.
The first axis is \emph{which} earlier blocks each block conditions on --- the sparsity that Vecchia,
GMRFs, banding and nested dissection all regularize along. The second is \emph{how strongly} that
conditioning is trusted, which is $\gamma$, and which those methods do not have: they run at full trust
and economize on structure. The two are orthogonal and compose.

\begin{center}\small
\begin{tabular}{lll}
\hline
 & $\gamma=1$ (full trust) & $\gamma=0$ (no trust)\\
\hline
dense conditioning & exact Gaussian likelihood; MVO $w\propto\Sigma^{-1}\mathbf 1$ & block composite likelihood; HRP\\
sparse conditioning & block-Vecchia / GMRF likelihood & (block-marginal, same column)\\
\hline
\multicolumn{3}{l}{\emph{interior $\gamma$: the Schur pseudo-likelihood --- the bridge none of the above carries}}\\
\hline
\end{tabular}
\end{center}

\noindent
The portfolio row is not an analogy. Minimum-variance optimization is $\gamma=1$ with full
conditioning; hierarchical risk parity is the $\gamma=0$ corner; Schur complementary allocation
interpolates them by damping the very same Schur complement with the very same $\gamma$. The Schur
pseudo-likelihood is that construction transported from risk to evaluation --- identical algebra, new
job --- and the reason an interior $\gamma$ beats the minimum-variance endpoint is the reason HRP is
steadier than raw mean--variance.

\section{It travels}

Nothing above needed finance, or even covariance estimation in particular. The ingredients are a
Gaussian (or Gaussian-pseudo) block likelihood and a partition, so the coupling-strength dial is
generic to Gaussian-block models. In spatial statistics and Gaussian processes, where Vecchia and
GMRFs already condition each site on a neighborhood, $\gamma$ adds a robustness knob on how far to
trust those neighbors, orthogonal to the neighborhood choice. In Gaussian graphical models it tempers
conditional-dependence strength between the full penalized likelihood and node-wise pseudo-likelihood.
In longitudinal and state-space models, with time slices as blocks, it interpolates between treating
slices independently and the full temporal joint. And anywhere composite likelihood is already used for
tractability, $\gamma$ is the dial from composite back toward full. I want to be clear that this last
paragraph is a claim about the mechanism, not a row of experiments; I have tested the
covariance-evaluation case and am pointing at the rest.

\section{Geometry, and what actually happened}

Damping interpolates each conditional covariance between $\Sigma_{kk}$ and $S_k$, and the obvious
straight line is not the only road. Moving instead along the affine-invariant SPD geodesic,
$\Sigma_{\mathrm{BD}}\,\#_\gamma\,\Sigma = \Sigma_{\mathrm{BD}}^{1/2}(\Sigma_{\mathrm{BD}}^{-1/2}
\Sigma\,\Sigma_{\mathrm{BD}}^{-1/2})^{\gamma}\Sigma_{\mathrm{BD}}^{1/2}$, lifts the small eigenvalues
multiplicatively rather than additively. The two agree at the ends and to first order, and part ways
near a singular Schur complement --- which is exactly the strong-coupling regime where it matters: the
linear damping shows a power dip there and the geodesic one does not.

As for results, the honest summary is three lines. The optimum in $\gamma$ moves with the regime: it
climbs to $1$ when the coupling is trustworthy, falls to $0$ when it is noise, and in between an
interior $\gamma$ strictly beats both ends, which is the bias--variance trade made visible. The
geodesic variant is the safer default where the two interpolations diverge. And on real equity returns
--- Ken French portfolios, rolling minimum-variance out-of-sample --- the estimators that invert a
raw covariance blow up by orders of magnitude while well-conditioned shrinkage stays put; that is not a
test of $\ell_\gamma$ as an equity estimator, and I will not pretend it is, but it is the same moral
the likelihood teaches, paid out in basis points: in high dimensions you are rewarded for not trusting
the raw inverse.

\section{What I can't yet claim}

The closed-form $\gamma^\star$ is real, but it is the \emph{estimator's} optimum --- the predictive,
well-behaved use. The high-dimensional \emph{discrimination} optimum is only half pinned down: the
variance term is first-order and I have it ($1/\lambda_{\min}^4$ at $\gamma=1$, bounded below it), but
the signal term is second-order in the sampling noise --- it is the shrinkage benefit --- so a closed
form for the discrimination $\gamma^\star$ wants the next order, and for now it is characterized only
numerically. The compounded $K$-block positive-definite threshold I have as a two-block formula and as
numerical evidence of compounding, not yet as a chain. Block ordering matters, as it does for Vecchia,
and I have not chased the permutation question. And the thesis is falsifiable in the right way: the
interior optimum should vanish whenever the coupling is fully reliable or pure noise, and an interior
optimum in a regime known to be fully reliable would sink it.

\section*{Acknowledgements}

The original Schur complementary theory, of which this is the evaluation-side mirror, was developed
with the support of Intech Investments.

\section*{References}

\small\noindent
L\'opez de Prado (2015, 2016); Antonov, Lipton \& L\'opez de Prado (2024); Cotton (2024,
arXiv:2411.05807); Besag (1975); Lindsay (1988); Varin, Reid \& Firth (2011); Vecchia (1988);
Katzfuss \& Guinness (2021); Pan et al.\ (2024); Rue \& Held (2005); George (1973); Ledoit \& Wolf
(2004, 2012); Bun, Bouchaud \& Potters (2017); James \& Stein (1961); Gneiting \& Raftery (2007);
Scheuerer \& Hamill (2015). Full annotated lists accompany the \texttt{precise} repository.

\end{document}
